\subsection{The exact remaining Selmer product for RH}
\label{shared-rh-selmer-product}

For the trivial target character the host can be written
\[
 E=E_h=\mathbf Q(\mu_{2^{h+2}})^+(\sqrt{-q},\sqrt{-d}),
 \quad \Delta=\operatorname{Gal}(E/\mathbf Q),\quad
 N=2^{h+2}qd,\quad [E:\mathbf Q]=4\cdot2^h.
\]
Here \(h\ge0\), and \(q,d\) are the distinct odd primes selected
by the arithmetic construction. Let \(P\) be its fixed finite set of
exceptional odd coefficient primes, containing \(3\), with
\(q,d>\max P\) and \(q,d>2t+1\), where \(t\ge2\) is the twist
cutoff. For \(p\) odd put \(S_{p,e}=\mathcal S(E,p,e)\),
\(1\le e<t\), and let \(j\in\Delta\) denote complex conjugation.

The strict condition can equivalently be imposed at every finite place.
Indeed, outside \(p\), the unramified local group is the quotient of
\(\mathbf Q_p/\mathbf Z_p(-e)\) by a nonzero Frobenius scalar minus
one, and this quotient is zero. We use the same full-local definition
over subfields when \(p\) is odd.

Since \(p\nmid|\Delta|\), the following are commuting integral
idempotents on \(S_{p,e}\):
\[
 C_e=\frac{1+(-1)^e j}{2},\qquad
 A_H=\frac1{|H|}\sum_{g\in H}g\quad(H\le\Delta).
\]
For \(p\in P\), define fixed thresholds and intermediate fields by
\[
 f_p=\begin{cases}1,&p\equiv1\pmod4,\\2,&p\equiv3\pmod4,\end{cases}
 \quad R_p=16p^2 4^{v_2(p^{f_p}-1)},\quad
 c_p=\max\{c\ge0:2^c<R_p\},
\]
\[
 F_p=\mathbf Q(\mu_{2^{\min(h+2,c_p)}})^+
                       (\sqrt{-q},\sqrt{-d}),
 \qquad H_p=\operatorname{Gal}(E/F_p).
\]
These fields retain the same \(q,d\); only their cyclotomic layers are
bounded independently of \(h\). Set
\begin{equation}\label{shared-rh-selmer-projectors}
 \Pi_{p,e}=\begin{cases}
 C_e,&p\notin P,\\
 C_e(1-A_\Delta),&p\in P,\ e\ge2,\\
 C_1(1-A_{H_p}),&p\in P,\ e=1,
 \end{cases}
 \qquad M_{p,e}=\Pi_{p,e}S_{p,e}.
\end{equation}
Every \(M_{p,e}\) is an integral direct factor of a finite group.

\begin{proposition}[Character coefficients and relative growth]
\label{shared-rh-integral-product}
Let \(\mathcal O\) be the integers of a finite extension of
\(\mathbf Q_p\) splitting \(\Delta\), with ramification index
\(b_{\mathcal O}\). For a character \(\theta\) of \(\Delta\), put
\(A_{\theta,e}=(\operatorname{Frac}\mathcal O/\mathcal O)
(\theta^{-1}\varepsilon_p^{-e})\), where \(\varepsilon_p\) is the
cyclotomic character, and let
\[
 \ell(p,e,\theta)=\frac1{b_{\mathcal O}}\operatorname{length}_{\mathcal O}
 \ker\left(H^1(\mathbf Q,A_{\theta,e})
       \longrightarrow\prod_{v<\infty}
         H^1(\mathbf Q_v,A_{\theta,e})\right).
\]
Then
\(\operatorname{length}_{\mathbf Z_p}M_{p,e}\) is the sum of
\(\ell(p,e,\theta)\) over the following characters:
\begin{enumerate}
\item \(p\notin P\): all \(\theta(-1)=(-1)^e\);
\item \(p\in P,e\ge2\): all nontrivial \(\theta(-1)=(-1)^e\);
\item \(p\in P,e=1\): the odd characters whose local conductor at
      \(2\) is at least \(R_p\).
\end{enumerate}
The equality is independent of the chosen splitting extension. Moreover,
\begin{equation}\label{shared-rh-relative-selmer}
 \#M_{p,1}=\frac{\#\mathcal S(E,p,1)^{j=-1}}
                  {\#\mathcal S(F_p,p,1)^{j=-1}}\quad(p\in P).
\end{equation}
\end{proposition}
\begin{proof}
The integral Fourier decomposition is an isomorphism because
\(|\Delta|\) is a unit in \(\mathcal O\). Its maps commute with
restriction to every decomposition group. The finite character
decomposition therefore identifies the kernel of the full localization
map with the direct sum of the twisted strict kernels, including the
ramified places. The projector \(C_e\) selects the stated parity,
\(A_\Delta\) the trivial character, and \(A_{H_p}\) the characters
factoring through \(F_p\). The quadratic factors have local conductor
at \(2\) at most \(4\), since their radicands are odd, and
\(c_p\ge2\). Thus these last characters are precisely those with
local conductor at \(2\) strictly less than \(R_p\).
Scalar extension multiplies integral lengths by
\(b_{\mathcal O}\), which cancels the displayed normalization.

For a Galois extension of degree prime to \(p\), inflation--restriction
identifies global first cohomology downstairs with the invariant part
upstairs: positive-degree cohomology of the finite quotient vanishes by
averaging. Local restriction is injective for the same reason, including
at ramified places. Thus restriction identifies the strict kernels.
Apply this to \(E/F_p\) and the commuting parity projector.
The invariant subgroup is a direct summand, so its order divides the
numerator in \eqref{shared-rh-relative-selmer}. If \(F_p=E\), the
quotient equals one.
\end{proof}

In the usual character partition, the only additionally omitted term
outside \(P\) is trivial, with \(p-1\mid2e\). Its removal uses the
separate strict trivial-character vanishing theorem for
\(p\ge5\), even \(e\ge2\), and \(p-1\mid2e\).
All its hypotheses hold: \(3\in P\) and trivial-character parity
forces \(e\) even. Thus \eqref{shared-rh-selmer-projectors} retains
exactly the three residual character families when that vanishing
theorem is supplied. It does not require that all trivial-character
groups vanish. A nontrivial odd-order character cannot occur because
\(\Delta\) is a \(2\)-group.

\subsection{The supplied remainder and the required allowance}
\label{shared-rh-selmer-allowance}

Keep the actual prime weights from
Proposition~\ref{shared-twisted-selmer-budget}, and put \(w_p=0\)
outside their finite support. Write
\[
 \mathcal V=\sum_p w_p\log p\sum_{e=1}^{t-1}
                  \operatorname{length}_{\mathbf Z_p}\mathcal S(E,p,e),
 \qquad
 \mathcal P(w)=\prod_{p\ne2}\prod_{e=1}^{t-1}(\#M_{p,e})^{w_p}.
\]
The weighted product is a positive real number; its exponents need not
be integers. The precise supplied-case hypotheses can be stated without
character coefficients. Let
\[
\begin{split}
 U_0={}&(\log2)\sum_{e<t}\operatorname{length}_{\mathbf Z_2}
                       \mathcal S(E,2,e)
       +\sum_{p\in P}\sum_{\substack{1\le e<t\\2\mid e}}
                            \log\#\mathcal S(\mathbf Q,p,e),\\
 U_-={}&\sum_{p\ne2}\sum_{e=1}^{t-1}\log\#((1-C_e)S_{p,e}),
 \qquad
 U_{\rm low}=\sum_{p\in P}\log\#\mathcal S(F_p,p,1)^{j=-1}.
\end{split}
\]
Here and below \(\sum_{e<t}\) means \(1\le e<t\).
Assume that the common-window, opposite-parity and bounded-layer
estimates, respectively, supply on the same whole interval
\begin{equation}\label{shared-rh-supplied-cases}
 U_0\le C_0e^{x^{4\varepsilon/5}}+C_1(t-1),\qquad
 U_-\le [E:\mathbf Q](t^2+t-2)\log(Nt),\qquad
 U_{\rm low}\le(Nt)^{C_s},
\end{equation}
with fixed nonnegative constants \(C_0,C_1,C_s\). These are explicit
inputs to the following assembly statement, not consequences of
qualitative finiteness.

\begin{proposition}[Exact residual budget]\label{shared-rh-residual-budget}
Under \eqref{shared-rh-supplied-cases},
\[
 0\le\mathcal V-\log\mathcal P(w)\le\mathcal K,
\]
where
\[
 \mathcal K=C_0e^{x^{4\varepsilon/5}}+C_1(t-1)
       +(Nt)^{C_s}+[E:\mathbf Q](t^2+t-2)\log(Nt).
\]
Fix \(0<\varepsilon<1/52\) and the constants in
\eqref{shared-rh-supplied-cases} before \(x\) varies. If
\(t\le x^{400}\), \([E:\mathbf Q]\le8x^{2+8\varepsilon}\), and
\(\log(Nt)\le2x^{\varepsilon/4}\), then
\(\mathcal K=o(e^{x^\varepsilon})\) as \(x\to\infty\), uniformly
on the stipulated windows. With \(y=x^\varepsilon\), \(B_0=fm e^y y\),
\(f,m\ge1\), \(N_c\ge0\), and a fixed \(C_V>0\), the sufficient
remaining condition is
\begin{equation}\label{shared-rh-product-target}
 \log\mathcal P(w)\le\tfrac12 C_V y(B_0+N_c).
\end{equation}
It implies \(\mathcal V\le C_Vy(B_0+N_c)\) eventually.
\end{proposition}
\begin{proof}
The complementary direct factors of \(M_{p,e}\) are exactly the
opposite-parity part, the rational invariant part for even \(e\) at
\(p\in P\), and the lower-layer odd part at \(e=1,p\in P\).
The prime \(2\) is retained separately. Orders multiply in each direct
sum. Removing the original weights, each between zero and one, bounds
the remainder by \(U_0+U_-+U_{\rm low}\), proving the first claim.

The opposite-parity contribution is at most
\(18t^2x^{2+33\varepsilon/4}\), since
\(t^2+t-2\le9t^2/8\). Also
\((Nt)^{C_s}\le\exp(2C_sx^{\varepsilon/4})\).
Every displayed contribution is \(o(e^{x^\varepsilon})\), uniformly
for \(t\le x^{400}\), with all constants fixed first.
Finally \(y(B_0+N_c)\ge e^y y^2\); the remainder is eventually
at most half the stated allowance.
\end{proof}

For a finite family sharing the same \(E,h,q,d,t\), only the prime
weights need change. If \(a_i\ge0\), \(\sum_i a_i=1\), and
\(\overline w_p=\sum_i a_iw_{i,p}\), then exactly
\[
 \sum_i a_i\log\mathcal P(w_i)=\log\mathcal P(\overline w).
\]
Actual rank factors must enter these averaging weights: for coefficients
\(u_i d_i\ge0\) with total \(D>0\), use
\(a_i=u_id_i/D\) and multiply the last equality by \(D\).
If \(D=0\), the weighted sum is zero before any division.
The same-host identification, the estimates
\eqref{shared-rh-supplied-cases}, the integral local comparisons and the
uniform bound \eqref{shared-rh-product-target} remain necessary.
Finiteness and this exact product identity do not supply the last bound.
